\documentclass[11pt]{article}
\usepackage{preamble}
\setlength{\parskip}{4pt}

\begin{document}

\daybanner{AP Statistics \textperiodcentered\ Unit 4 \textperiodcentered\ Topic 4.10 \textperiodcentered\ B: 2027-01-29 \textperiodcentered\ E: 2027-03-19 \textperiodcentered\ TEACHER KEY}{What Does This Test Support?}

\sectionbanner{CED TETHER}

{\small
\begin{itemize}[leftmargin=1.5em,itemsep=2pt,topsep=2pt]
  \item \textbf{Skill 3.E} --- Calculate a test statistic and find a p-value, provided conditions for inference are met.
  \item \textbf{Skill 4.B} --- Interpret statistical calculations and findings to assign meaning or assess a claim.
  \item \textbf{Skill 4.E} --- Justify a claim using a decision based on significance tests.
  \item \textbf{EU VAR-7} --- The t-distribution may be used to model variation.
  \item \textbf{LO VAR-7.I} --- Calculate an appropriate test statistic for a difference of two means. [Skill 3.E]
  \item \textbf{EK VAR-7.I.1} --- For a single quantitative variable, data collected using independent random samples or a randomized experiment from two populations, each of which can be modeled with a normal distribution, the sampling distribution of t = [($\bar{x}_1$ - $\bar{x}_2$) - ($\mu_1$ - $\mu_2$)] / $\surd$($s_1^2$/$n_1$ + $s_2^2$/$n_2$) is an approximate t-distribution with degrees of freedom that can be found using technology. The degrees of freedom fall between the smaller of $n_1$ - 1 and $n_2$ - 1 and $n_1$ + $n_2$ - 2.
  \item \textbf{EU DAT-3} --- Significance testing allows us to make decisions about hypotheses within a particular context.
  \item \textbf{LO DAT-3.G} --- Interpret the p-value of a significance test for a difference of population means. [Skill 4.B]
  \item \textbf{EK DAT-3.G.1} --- An interpretation of the p-value of a significance test for a two-sample difference of population means should recognize that the p-value is computed by assuming that the null hypothesis is true, i.e., by assuming that the true population means are equal to each other.
  \item \textbf{LO DAT-3.H} --- Justify a claim about the population based on the results of a significance test for a difference of two population means in context. [Skill 4.E]
  \item \textbf{EK DAT-3.H.1} --- A formal decision explicitly compares the p-value to the significance $\alpha$. If the p-value $\leq$ $\alpha$, then reject the null hypothesis, $H_0$: $\mu_1$ - $\mu_2$ = 0, or $H_0$: $\mu_1$ = $\mu_2$. If the p-value > $\alpha$, then fail to reject the null hypothesis.
  \item \textbf{EK DAT-3.H.2} --- The results of a significance test for a two-sample test for a difference between two population means can serve as the statistical reasoning to support the answer to a research question about the populations that were sampled.
\end{itemize}
}
\textbf{Essential question:} How can one test support a mean difference without settling size, individuals, or cause?\par
\textbf{DOK-3 skill:} 4.E \textperiodcentered\ \textbf{FRQ pattern:} {\small\texttt{significance-\allowbreak{}vs-\allowbreak{}size-\allowbreak{}cause-\allowbreak{}and-\allowbreak{}scope}}\par
\textbf{DOK rationale:} Part (c) synthesizes a warranted report from two incomplete accounts by coordinating the p-value decision, practical benchmark, mean-versus-individual scope, sampling, and self-selection.\par

\frameworkphaseheader{Do Now (first take) $\rightarrow$ Explore (video + follow-along) $\rightarrow$ Exit (finish + turn in)}{1 $\rightarrow$ 3}{45}{%
\begin{itemize}[leftmargin=1.5em,itemsep=2pt,topsep=2pt]
  \item Display the board and ask for one sentence using the study description.
  \item Begin the 7.9 video follow-along at minute 5.
  \item Return to the board at minute 33. Students finish the shortened parts and justify their judgment.
  \item Collect at the door. Score part (c) only using the three required elements.
\end{itemize}
}{%
\begin{itemize}[leftmargin=1.5em,itemsep=2pt,topsep=2pt]
  \item Write a one-sentence first take before the video.
  \item Complete the video follow-along worksheet.
  \item Finish (a)--(c), revise the first take in the exit line, and turn in the sheet.
\end{itemize}
}{%
\begin{itemize}[leftmargin=1.5em,itemsep=2pt,topsep=2pt]
  \item Which specific number or study detail supports your choice?
  \item What claim would go beyond the evidence, and why?
\end{itemize}
}{Circulate and ask for evidence. Accept a concise response; do not require omitted calculations or a second full inference write-up.}

\begin{callout}{\IconWarn\ WATCH FOR}{calloutred}
\begin{itemize}[leftmargin=1.5em,itemsep=2pt,topsep=2pt]
  \item Choosing a speaker without a reason.
  \item Copying numbers without explaining what they support.
\end{itemize}
\end{callout}

\sectionbanner{SCORING PART (c) --- E / P / I}

\textbf{Expected elements}
\begin{itemize}[leftmargin=1.5em,itemsep=2pt,topsep=2pt]
  \item \textbf{decision} (required) --- Uses p < 0.05 to retain Nia's population-mean conclusion
  \item \textbf{size} (required) --- Retains Mateo's limit using 1.2 versus 3 and the test against zero
  \item \textbf{cause-individuals} (required) --- Explains why a mean difference and self-selection do not justify a universal causal claim
\end{itemize}

{\small\renewcommand{\arraystretch}{1.25}
\noindent\begin{tabular}{|p{0.5in}|p{5.9in}|}
\hline
\textbf{E} & Combines the correct clauses using the p-value, practical benchmark, and the limits on individual and causal claims. \\ \hline
\textbf{P} & Shows substantive valid reasoning for at least one required element, but does not meet all E requirements. Credit correct reasoning even when another claim is wrong. \\ \hline
\textbf{I} & Shows no substantive valid reasoning for any required element; a name, unsupported choice, or copied number alone does not count. \\ \hline
\end{tabular}}

\textbf{Common mistakes}
\begin{itemize}[leftmargin=1.5em,itemsep=2pt,topsep=2pt]
  \item Reading 0.0263 as $P(H_0)$
  \item Treating a test against 0 as a test against 3
  \item Inferring cause from mode choice
\end{itemize}

\clearpage
\focusbanner{TODAY'S PROBLEM --- annotated}

Suppose this term 4{,}000 students chose Focus mode and 3{,}600 chose Guided. Independent random samples of 100 students from each mode completed the same practice set. The service asked in advance whether mean Focus time was greater and set $\alpha=0.05$. Only a difference of at least 3 minutes matters in practice.\par\smallskip \textbf{Nia:} ``The $p\approx0.0263<0.05$ supports a higher Focus mean. So Guided causes every student to finish at least 3 minutes faster; switch everyone.''\par \textbf{Mateo:} ``The sample gap is 1.2 minutes, and students chose their modes. A 3-minute effect and causation are not established. So retain $H_0$; there is no evidence of different population means.''\par
\begin{center}
\textbf{Practice time summaries (minutes)}\par\smallskip
\begin{tabular}{|l|c|c|c|c|}
\hline
\textbf{Mode} & \textbf{$N$} & \textbf{$n$} & \textbf{$\bar x$} & \textbf{$s$} \\ \hline
\textbf{Focus} & 4{,}000 & 100 & 42.8 & 4.2 \\ \hline
\textbf{Guided} & 3{,}600 & 100 & 41.6 & 4.5 \\ \hline
\end{tabular}
\end{center}
\begin{firsttakebox}{FIRST TAKE (5 min)}
Which claim seems too strong: no average difference, or Guided helps everyone? Give one reason.\par
\answer{Ungraded. Everyday language is enough before the video; formal vocabulary belongs in the finish.}
\end{firsttakebox}

\textit{Rules callout ``CARRY OUT THE TEST, THEN BOUND THE CONCLUSION (keep on the page)'' is printed on the student sheet.}\par\medskip

\dokbadge{1}~\textbf{(a)} State $H_0$ and $H_a$ for the population mean times in Focus-minus-Guided order.\par
\answer{$H_0:\mu_F-\mu_G=0$; $H_a:\mu_F-\mu_G>0$, for the mean times of all active Focus and Guided choosers this term.}

\dokbadge{2}~\textbf{(b)} The difference is 1.2 minutes and $SE=0.6155485$. Calculate $t$ and interpret the supplied one-sided $p=0.0263$ in one sentence.\par
\answer{$t=1.2/0.6155485=1.9495$. Assuming equal population means, the chance of a test statistic at least this large is about 0.0263.}

\dokbadge{3}~\textbf{(c)} In 4 sentences, combine the supported parts of both reports. Use the p-value for the mean claim, compare 1.2 with the 3-minute benchmark, and explain why self-selection and a mean difference do not justify switching every student.\par
\answer{Nia correctly uses $0.0263<0.05$: reject $H_0$ and report evidence of a higher mean time for this term's Focus choosers. Mateo correctly notes that the observed gap is only 1.2 minutes and the test against 0 does not establish a 3-minute difference. A population mean difference does not show that every Guided student is faster. Because students selected their modes, the comparison does not establish that switching modes causes improvement, so neither report should be kept unchanged.}

\begin{summaryexitbox}{EXIT}
One thing the video changed about my first take: \hrulefill
\end{summaryexitbox}

\end{document}
